% Complete product-measure and iterated-integration proofs for P15-U1.
\section{Product measures and iterated integration: proofs}
\label{sec:annex-p15-product-measures}

Use Chapter~\ref{sec:p15-product-measures} for the product notation.
The construction uses only the earlier measure-extension theorem and
one-variable nonnegative integration. In particular, its proof precedes
and does not assume Tonelli or Fubini. The source is Durrett's fixed
11 January 2019 author version \cite[sec.~1.7]{DurrettPTEAuthor2019}.

\begin{annexitem}{annex:p15.product-sections}{thm:p15-product-sections}{Sections, coordinate maps and Borel products}
\begin{proof}
Fix \(x\in X\). The collection of subsets \(E\subseteq X\times Y\)
for which \(E_x\in\mathcal B\) is a sigma-algebra: the whole
product has section \(Y\), and
\[
 (E^c)_x=Y\setminus E_x,\qquad
 \left(\bigcup_{n\geq1}E_n\right)_x=\bigcup_{n\geq1}(E_n)_x.
\]
It contains every rectangle, whose section is \(B\) when \(x\in A\)
and empty otherwise. It therefore contains \(\mathcal A\times\mathcal B\).
This works separately for every \(x\), without discarding any point.
The same argument with a fixed \(y\) proves the other section assertion.

For a product measurable extended real \(f\) and any real \(a\),
\[
 \{y:f_x(y)\leq a\}=\{(u,v):f(u,v)\leq a\}_x\in\mathcal B.
\]
The extended-real measurability criterion proves that \(f_x\) is
measurable; the proof for \(f^y\) is identical with the variables
interchanged. Real-valued functions use the real version of the criterion.

The inverse images of \(A\in\mathcal A\) and \(B\in\mathcal B\)
under the two projections are respectively \(A\times Y\) and
\(X\times B\). These rectangles are measurable, so the projections
are measurable. Conversely, any sigma-algebra making them measurable
contains their inverse-image intersections \(A\times B\), hence
contains the generated product sigma-algebra. If a map \(T\) is
measurable, composing with either projection preserves measurability.
If both coordinate maps are measurable, then
\[
 T^{-1}(A\times B)
   =(\pi_X\circ T)^{-1}(A)\cap(\pi_Y\circ T)^{-1}(B)
\]
is measurable for every rectangle. The generator criterion from
Theorem~\ref{thm:p15-measurability-criterion} proves measurability of \(T\).

On \(\mathbb R^{r+s}\), coordinate-block projections are continuous
and therefore Borel measurable. Each rectangle with Borel factors is
the intersection of two Borel inverse images, which proves that the
product sigma-algebra is contained in \(\mathcal B(\mathbb R^{r+s})\).
Conversely, every open rectangle with rational endpoints in
\(\mathbb R^{r+s}\) is a product of rational open rectangles in
the two coordinate blocks. These sets generate the Euclidean Borel
sigma-algebra by Theorem~\ref{thm:p15-borel-generators}, proving the
reverse containment. No measure was used in these arguments.
\end{proof}
\end{annexitem}

\begin{annexitem}{annex:p15.product-existence}{thm:p15-product-existence}{The rectangle premeasure and its extension}
\begin{proof}
The rectangles contain the empty set and are closed under intersections:
\((A\times B)\cap(C\times D)=(A\cap C)\times(B\cap D)\).
The complement of a rectangle is the disjoint union
\[
 (A^c\times B)\ \cup\ (A\times B^c)\ \cup\ (A^c\times B^c),
\]
where complements are in their respective factors. Hence the rectangles
form a semialgebra. A nonempty rectangle determines both factors as its
coordinate projections, so its assigned value is independent of its
description. An empty rectangle has at least one empty factor and its
value is zero under the prescribed convention. This proves that \(m\)
is well defined and \(m(\varnothing)=0\).

We first record the one-variable identity
\[
 \int c\mathbf1_C\,d\mu_i=c\mu_i(C)
       \qquad(c\in[0,\infty],\ C\text{ measurable}).
\]
For finite \(c\), this is the simple integral formula. For
\(c=\infty\), the functions \(k\mathbf1_C\) increase to the
specified extended function. Monotone convergence gives zero when
\(\mu_i(C)=0\) and infinity otherwise, which is precisely the
right side. Thus this identity needs no unproved rule for an infinite
coefficient in a simple function.

Suppose a rectangle is a finite or countable disjoint union of
rectangles, padded by empty sets in the finite case:
\[
 A\times B=\bigcup_{j\geq1}(A_j\times B_j).
\]
At a fixed \(x\), the sections on the right are pairwise disjoint
and have union the section on the left. Countable additivity in \(Y\)
therefore gives the pointwise extended identity
\[
 \mathbf1_A(x)\mu_2(B)
        =\sum_{j\geq1}\mathbf1_{A_j}(x)\mu_2(B_j).
\]
This is valid even outside \(A\): any nonempty section on the right
there would contradict the assumed set equality. All functions in this
identity are nonnegative and measurable on \(X\). Integrating, and
using the proved nonnegative series law
(Theorem~\ref{thm:p15-nonnegative-integral-laws}) and the constant
indicator identity, gives
\[
 \mu_1(A)\mu_2(B)
       =\sum_{j\geq1}\mu_1(A_j)\mu_2(B_j).
\]
Thus both semialgebra partition conditions hold, including infinite
values and zero factors. Theorem~\ref{thm:p15-semialgebra-extension}
gives a premeasure on the algebra of finite disjoint unions of rectangles.
The countable-cover construction in
Theorem~\ref{thm:p15-premeasure-extension} extends it to
\(\sigma(\mathcal S)=\mathcal A\times\mathcal B\). That theorem's
existence branch requires no sigma-finiteness, proving the first assertion.

Now assume the factors are sigma-finite. Choose increasing measurable
sets \(F_n\uparrow X\), \(G_n\uparrow Y\) with
\(\mu_1(F_n)<\infty\), \(\mu_2(G_n)<\infty\). Such increasing
covers are obtained from any sigma-finite covers by finite unions.
The rectangles \(F_n\times G_n\) increase to \(X\times Y\):
both coordinates of a fixed point have entered their respective covers
by some common index. Their assigned masses are finite products of
finite numbers. Thus the rectangle-algebra premeasure is sigma-finite,
so its extension is unique by the uniqueness branch of the extension
theorem. Any measure with the required rectangle values agrees on
finite disjoint rectangle unions and is such an extension. These same
rectangles prove sigma-finiteness of \(\rho\). Finally,
\(\rho(X\times Y)=\mu_1(X)\mu_2(Y)=1\) for probability factors.
\end{proof}
\end{annexitem}

\begin{annexitem}{annex:p15.finite-product-measures}{thm:p15-finite-product-measures}{Construction, uniqueness and coordinate identifications}
\begin{proof}
For \(n=1\), the assertion is the given measure space. Suppose the
product of the first \(n-1\) factors has been constructed and is
sigma-finite. Form its binary product with the last factor using the
preceding theorem. Its sigma-algebra is the one generated by all
\(n\)-fold rectangles, as we now check.

Let \(\mathcal H\) be the sigma-algebra generated by those rectangles.
For fixed \(A_n\in\mathcal F_n\), the collection of subsets \(C\)
of the first \(n-1\) factors such that \(C\times A_n\in\mathcal H\)
is a sigma-algebra: the whole-set rectangle belongs to \(\mathcal H\),
relative complementation follows by subtracting \(C\times A_n\)
from that rectangle, and countable unions commute with multiplication
by \(A_n\). It contains every \((n-1)\)-fold rectangle, and hence
their generated sigma-algebra. Consequently every binary rectangle
for the grouped construction belongs to \(\mathcal H\). The reverse
containment is immediate because every \(n\)-fold rectangle is a
binary rectangle for that grouping. This proves equality of domains.
The binary rectangle formula and the induction hypothesis prove the
\(n\)-fold mass formula, and the binary theorem proves sigma-finiteness.

For uniqueness at all \(n\), choose for each factor an increasing
finite-measure cover \(F_{i,k}\), \(k\geq1\). The rectangles
\(H_k=\prod_{i=1}^n F_{i,k}\) increase to the whole product and have
finite prescribed masses. If two measures agree on the rectangles,
fix \(k\) and consider the sets \(E\) for which their values on
\(E\cap H_k\) agree. The whole space is in this class. Nested
differences preserve equality by subtracting finite values, and increasing
unions preserve it by continuity from below. The class is therefore a
lambda-system containing the generating rectangle pi-system. The
pi-lambda theorem makes it contain every product measurable \(E\).
Letting \(k\) increase then proves equality of the measures on \(E\),
again by continuity from below.

For a division into two blocks of coordinates, fix first an elementary
rectangle in the second block. The preceding sigma-algebra argument
extends the first block's rectangles to all its measurable sets.
Now fix any such measurable first-block set. The same argument in
the second block extends its rectangles to all its measurable sets;
the whole second block was already an elementary rectangle, so the
relative-complement step is justified. Thus every grouped binary
rectangle belongs to the original generated domain, and the reverse
containment follows from the elementary rectangles. Any binary
parenthesization therefore has the same domain. Induction on the
sizes of its two blocks gives the same rectangle masses. Uniqueness
therefore identifies the resulting measures. A coordinate permutation
and its inverse carry rectangles to rectangles and are measurable.
Transporting the measure by that bijection preserves countable additivity
and gives the prescribed permuted rectangle masses. Uniqueness gives the
permutation assertion. The nonnegative extended finite product is
independent of grouping: if any factor is zero the whole product is zero;
otherwise it is the ordinary positive product or infinity. Thus its
conventions also agree under these identifications.

For Euclidean factors, successive applications of the Borel-product
identity give the domain \(\mathcal B(\mathbb R^n)\). Each
\(\lambda\) is the already constructed sigma-finite Borel measure,
and its interval values are \(b_i-a_i\). Substitution in the
rectangle formula proves the displayed volume identity. This constructs
the stated Borel Lebesgue measure; it makes no assertion about completing
its domain.
\end{proof}
\end{annexitem}

\begin{annexitem}{annex:p15.section-measures}{thm:p15-section-measures}{Finite rectangles, the pi-lambda argument and exhaustion}
\begin{proof}
Choose increasing finite-measure covers \(F_n,G_n\) as in the product
existence proof, and put \(H_n=F_n\times G_n\). Fix \(n\).
For product measurable \(E\), define
\[
 g_{n,E}(x)\defeq
 \begin{cases}
  \mu_2(E_x\cap G_n),&x\in F_n,\\
  0,&x\notin F_n.
 \end{cases}
\]
Every section in this formula is measurable by the section theorem,
so its measure is defined. Let \(\mathcal L_n\) consist of those
\(E\in\mathcal A\times\mathcal B\) for which \(g_{n,E}\)
is \(\mathcal A\)-measurable and
\(\int_X g_{n,E}\,d\mu_1=\rho(E\cap H_n)\).

For a rectangle \(E=A\times B\), the function is
\(\mu_2(B\cap G_n)\mathbf1_{A\cap F_n}\). It is measurable,
and its integral is
\(\mu_1(A\cap F_n)\mu_2(B\cap G_n)=\rho(E\cap H_n)\).
Thus the rectangle pi-system belongs to \(\mathcal L_n\), including
the whole product. If \(E\subseteq D\) are members, their sections
inside \(G_n\) are nested and of finite measure. Consequently
\[
 g_{n,D\setminus E}=g_{n,D}-g_{n,E}.
\]
The functions on the right are real and integrable: they are bounded
by the finite constant \(\mu_2(G_n)\) on the finite-measure set
\(F_n\) and vanish outside it. Signed linearity therefore applies,
and finite-measure subtraction in \(H_n\) gives
\[
 \int g_{n,D\setminus E}\,d\mu_1
  =\rho(D\cap H_n)-\rho(E\cap H_n)
  =\rho((D\setminus E)\cap H_n).
\]
This proves closure under nested differences without ever subtracting
infinite values.

If \(E_j\in\mathcal L_n\) increases to \(E\), continuity from
below in the second factor gives \(g_{n,E_j}\uparrow g_{n,E}\)
pointwise. Measurable limits and monotone convergence imply that the
limit function is measurable and its integral is
\[
 \lim_j\rho(E_j\cap H_n)=\rho(E\cap H_n).
\]
Thus \(\mathcal L_n\) is a lambda-system. By the already proved
pi-lambda theorem it contains \(\mathcal A\times\mathcal B\).

Now fix arbitrary product measurable \(E\) and let \(n\) increase.
At every \(x\), eventual membership in \(F_n\) and continuity
from below on \(E_x\cap G_n\) show
\(g_{n,E}(x)\uparrow\mu_2(E_x)\). The limit is measurable, and
monotone convergence, followed by continuity from below in the product,
gives
\[
 \int_X\mu_2(E_x)\,d\mu_1(x)
   =\lim_n\rho(E\cap H_n)=\rho(E).
\]
Interchanging the two factors gives the other equality; their product
measures agree under coordinate exchange by the finite-product theorem.
Finally, the nonnegative zero-integral criterion gives
\(\int_X\mu_2(E_x)\,d\mu_1=0\) exactly when
\(\mu_2(E_x)=0\) for \(\mu_1\)-almost every \(x\).
Applying the same criterion in the other order gives the last assertion.
\end{proof}
\end{annexitem}

\begin{annexitem}{annex:p15.tonelli}{thm:p15-tonelli}{Nonnegative approximation in the product and both factors}
\begin{proof}
For \(f=\mathbf1_E\), the two section integrals are the section
measure functions of the preceding lemma. That lemma proves their
measurability and both integral equalities. For a nonnegative simple
\(f=\sum_{j=1}^m a_j\mathbf1_{E_j}\), finite nonnegative additivity
and scaling in each factor and in the product prove the same statements.
The coefficients \(a_j\) are finite, and all sums use nonnegative
extended arithmetic, so this step permits infinite measure regions.

For a general nonnegative product measurable \(f\), take its explicit
increasing nonnegative simple approximations \(s_n\uparrow f\) from
Theorem~\ref{thm:p15-measurable-limits}, applied on \(X\times Y\).
For every fixed \(x\), the sections \((s_n)_x\) increase to
\(f_x\); they are measurable, as is \(f_x\), by the section
theorem. One-variable monotone convergence gives
\[
 T_1s_n(x)\uparrow T_1f(x).
\]
Each function on the left is measurable by the simple case, so the
limit is measurable. Monotone convergence now in \(X\), the simple
case, and monotone convergence in the product give
\[
 \int_X T_1f\,d\mu_1
   =\lim_n\int_X T_1s_n\,d\mu_1
   =\lim_n\int s_n\,d\rho
   =\int f\,d\rho.
\]
Applying the same three steps with \(Y\) outside gives the second
equality. Each use concerns nonnegative functions and permits infinity.

For the last assertion, let \(f,g\) agree outside a measurable product
null set \(N\). The section-measure lemma makes \(N_x\) null for
almost every \(x\), and \(N^y\) null for almost every \(y\).
Outside the first exceptional set the two measurable functions
\(f_x,g_x\) agree a.e. and therefore have equal nonnegative integrals
by Theorem~\ref{thm:p15-nonnegative-integral-laws}. The same argument
in the other factor proves the second a.e. equality. No claim about
the section integrals at every exceptional point follows from this argument.
\end{proof}
\end{annexitem}

\begin{annexitem}{annex:p15.fubini}{thm:p15-fubini}{The signed theorem and its measurable exceptional sets}
\begin{proof}
Apply Tonelli to \(|f|\). It makes \(H_1,H_2\) measurable and gives
\[
 \int_XH_1\,d\mu_1=\int|f|\,d\rho=\int_YH_2\,d\mu_2<\infty.
\]
The finite-integral criterion in
Theorem~\ref{thm:p15-nonnegative-integral-laws} proves that the
measurable sets \(N_1,N_2\) are null. Outside them the section
absolute integral is finite, which is exactly section integrability.
The section theorem also applies at points inside those null sets;
it is integrability, not measurability, that may fail there.

Tonelli applies separately to \(f^+\) and \(f^-\), producing
nonnegative measurable functions
\(P_1=T_1(f^+)\), \(Q_1=T_1(f^-)\).
Their sum is \(H_1\) pointwise, by nonnegative additivity in each
section. Hence both are finite on \(N_1^c\). Replace both by zero
on \(N_1\), calling the resulting functions \(\widetilde P_1\),
\(\widetilde Q_1\). They are measurable real functions: this is a
piecewise definition on a measurable set, and finite values occur on
each piece. Their integrals equal those of \(P_1,Q_1\), because
the replacements occur on a measurable null set. Tonelli gives
\[
 \int_X\widetilde P_1\,d\mu_1=\int f^+\,d\rho<\infty,
 \qquad
 \int_X\widetilde Q_1\,d\mu_1=\int f^-\,d\rho<\infty.
\]
The function \(I_1f\) in the statement is precisely the everywhere
finite difference \(\widetilde P_1-\widetilde Q_1\). Signed
linearity and the definition of the signed product integral yield
\[
 \int_X I_1f\,d\mu_1
   =\int f^+\,d\rho-\int f^-\,d\rho
   =\int f\,d\rho.
\]
All the subtracted numbers are finite. The bound
\(|I_1f|\leq\widetilde P_1+\widetilde Q_1\) also gives
\(\int_X|I_1f|\,d\mu_1\leq\int|f|\,d\rho\).
The same construction with \(P_2=T_2(f^+)\), \(Q_2=T_2(f^-)\)
and \(N_2\) proves every assertion in the other order. This proof
includes integrable extended real \(f\); it never forms an undefined
pointwise difference of infinite section integrals.

If another product measurable function \(g\) equals \(f\) outside
a measurable product null set, its absolute integral is the same and
finite, so the theorem applies to it as well. The section-measure
lemma makes their sections equal a.e. for almost every point in each
factor. Remove also the two null sets on which either section absolute
integral is infinite. On the remaining full-measure sets, signed a.e.
invariance of integrals proves equality of the respective \(I_i\).
Finally, any other measurable finite assignment on the specified
\(N_i\) differs from \(I_i f\) only on a measurable null set.
It is therefore integrable and has the same integral by signed a.e.
invariance. Neither step requires completeness of a factor measure.
\end{proof}
\end{annexitem}

\begin{annexitem}{annex:p15.fubini-failures}{thm:p15-fubini-failures}{The four failures and the elementary integral calculations}
\begin{proof}
First justify the one-dimensional calculations used here and in
Example~\ref{def:p15-triangle-integral}. For a continuously
differentiable \(h\) on a compact interval \([a,b]\), let
\(L\) bound \(|h'|\). On a partition into \(m\) equal intervals,
the mean value theorem bounds the oscillation on each cell by
\(L(b-a)/m\). Lower and upper simple functions therefore bracket
\(h\) and have integral difference at most \(L(b-a)^2/m\).
Their finite weighted sums converge to the ordinary Riemann integral,
so the Lebesgue integral has that value. Endpoint singletons have
zero measure. The fundamental theorem of elementary calculus now
evaluates these integrals by an antiderivative. Thus, for example,
\[
 \int_a^b t\,d\lambda(t)=(b^2-a^2)/2,\qquad
 \int_a^b t^2\,d\lambda(t)=(b^3-a^3)/3,
\]
and for \(c>0\),
\(\int_a^b e^{-ct}\,d\lambda(t)=(e^{-ca}-e^{-cb})/c\).
The exponential formula extends to \(b=\infty\) by monotone
convergence on increasing bounded intervals. The derivatives and
fundamental theorem used in these numerical examples are ordinary
one-variable calculus; the product-measure proofs above do not depend
on them. The polynomial formulas give all steps of the unit-square
calculation in either order.

For the first counterexample, the product sigma-algebra on
\(\mathbb N_+^2\) contains every singleton and thus every subset,
since the set of ordered pairs is countable. The product of the two
counting measures is counting measure on the pairs: singleton masses
are one, and countable additivity determines all masses. The function
\(f\) is consequently measurable. For \(m=1\), the section
\(n\mapsto f(m,n)\) has sum one. For \(m\geq2\), it has
values \(-1\) at \(n=m-1\) and \(1\) at \(n=m\), so
its sum is zero. The first iterated sum is therefore one. For each
\(n\), the other section has \(1\) at \(m=n\) and \(-1\)
at \(m=n+1\), hence sum zero. Its outer sum is zero. Each
positive diagonal point contributes one to \(\int f^+\), and
each negative point contributes one to \(\int f^-\); there are
infinitely many of both. Both nonnegative product integrals are
infinite, so the signed product integral is undefined.

For the second example, \(f\) is continuous on its product of open
intervals, and is product measurable by the Borel-product theorem.
For fixed \(x>0\), both positive exponential terms are integrable
over \((1,\infty)\), hence their difference is integrable. Their
elementary integrals give
\[
 \int_1^\infty f(x,y)\,d\lambda(y)
                  =(e^{-x}-e^{-2x})/x.
\]
This is positive for \(0<x<1\) and bounded above by one, since
\(e^{-x}(1-e^{-x})\leq x\). The last inequality follows by
integrating \(e^{-t}\leq1\) from zero to \(x\).
Thus the first outer integral is finite and strictly positive; strict
positivity also follows from the nonnegative zero-integral criterion
because the integrand is positive throughout an interval of positive measure.
For fixed \(y>1\), both exponentials are bounded and integrable
over \((0,1)\), and subtraction of their elementary integrals gives
\[
 \int_0^1 f(x,y)\,d\lambda(x)
                  =(e^{-2y}-e^{-y})/y.
\]
This is strictly negative and its absolute value is at most \(e^{-y}\),
whose integral on \((1,\infty)\) is finite. The second outer integral
is therefore finite and strictly negative.

To identify the product signed-integral failure directly, put
\(\phi(t)=e^{-t}-2e^{-2t}\). This continuous function is positive
on \([1,2]\) and negative on \([1/4,1/2]\): its sign is that of
\(1-2e^{-t}\), with change at \(\log 2\), which lies between
\(1/2\) and \(1\). Let \(c_+>0\) and \(c_->0\) be respective
lower bounds for \(\phi\) on the first compact interval and
\(-\phi\) on the second. Such bounds exist by continuity and strict
sign on those intervals. For \(0<x<1\), the interval
\([1/x,2/x]\) lies in \((1,\infty)\) up to an irrelevant endpoint
and has length \(1/x\). On it \(f(x,y)\geq c_+\), so
\(\int_1^\infty f^+(x,y)\,d\lambda(y)\geq c_+/x\).
For \(0<x<1/4\), the interval \([1/(4x),1/(2x)]\) lies in
the same \(y\)-domain, has length \(1/(4x)\), and satisfies
\(f^-(x,y)\geq c_-\). Its section integral is at least
\(c_-/(4x)\). The integral of \(1/x\) on \((0,a)\), for
any \(a>0\), is infinite: on each disjoint interval
\((a2^{-(k+1)},a2^{-k}]\), its integral is at least \(1/2\),
by bounding it below by \(2^k/a\). Countable nonnegative
additivity proves divergence. Tonelli applied separately to \(f^+\)
and \(f^-\) now proves that both product part integrals are infinite.

For the third example, counting measure on all subsets of an
uncountable set is a measure: a disjoint union with infinite cardinality
either contains an infinite member or has arbitrarily large finite
partial cardinality sums, and finite unions have their ordinary
cardinality sums. A set of finite counting measure is finite, so a
countable union of such sets is countable and cannot cover \((0,1)\).
This factor is not sigma-finite. The first factor is a probability
measure. The diagonal is product measurable: its complement is
\[
 \{x<y\}\cup\{y<x\},\qquad
 \{x<y\}=\bigcup_{q\in\mathbb Q}
     \bigl((0,1)\cap(-\infty,q)\bigr)
       \times\bigl((0,1)\cap(q,\infty)\bigr),
\]
and the analogous rational-rectangle union describes \(\{y<x\}\).
All the first-factor sets are Borel and the second-factor sets are
measurable in their full power set. A vertical section of the diagonal
is the singleton \(\{x\}\), of counting measure one. A horizontal
section is the singleton \(\{y\}\), of Lebesgue measure zero.
The two iterated values are therefore respectively one and zero.
They are well-defined nonnegative integrals, but cannot both equal
one common value of any product-space integral. Existence of the
specified rectangle extension from the earlier theorem does not
rescue the exchange equality when the sigma-finiteness premise fails.

For the last example, keep the explicitly assumed bijection \(b\).
If \(y=b(\beta)\), the set \(\{x:x<'y\}=b[\beta]\) is
countable because \(\beta\) is a countable ordinal. It is a
countable union of Borel singletons and has measure zero. If
\(x=b(\alpha)\), the complement in \((0,1)\) of
\(\{y:x<'y\}\) is \(b[\alpha+1]\), also countable and null.
Thus that section is Borel with measure one. These prove all the
asserted section integrals and make both remaining-variable functions
constant and measurable. If the joint indicator were product measurable,
Tonelli for these probability factors would make the two iterated
integrals equal, contradicting zero versus one. It is therefore not
product measurable. The argument uses the stated set-theoretic
assumptions exactly to obtain the bijection and the countable initial
segments; it does not remove them from the example's premise.
\end{proof}
\end{annexitem}

\begin{annexitem}{annex:p15.absolute-iterated-series}{thm:p15-absolute-iterated-series}{Testing absolute integrability and exchanging a series}
\begin{proof}
Tonelli applied to the nonnegative product measurable function \(|f|\)
identifies both iterated absolute integrals with its product integral.
Thus the three finiteness conditions are equivalent. Any one of them
supplies exactly the absolute-integrability premise of Fubini, whose
construction also supplies measurable versions of the signed section
integrals and their necessary a.e. qualification.

For the series assertion, put counting measure \(\kappa\) on
\(\mathbb N_+\), and define \(F(n,y)=f_n(y)\).
For every real \(a\),
\[
 \{(n,y):F(n,y)\leq a\}
       =\bigcup_{n\geq1}\{n\}\times\{y:f_n(y)\leq a\},
\]
so \(F\) is product measurable. Both factors are sigma-finite.
The assumed sum is the first iterated absolute integral, because
integration with respect to \(\kappa\) is summation. It is finite,
so the preceding criterion and Fubini apply. The function
\(H(y)=\sum_n|f_n(y)|\) is measurable by nonnegative measurable
limits and has finite integral by Tonelli. Hence
\(N=\{H=\infty\}\) is measurable and null. For \(y\notin N\),
ordinary absolute convergence defines \(\sum_n f_n(y)\), and
counting-measure integration of that integrable section equals the sum.
Fubini's other section-integral function, with zero assigned on \(N\),
is exactly the stated \(S\). In particular it is measurable and
integrable, and Fubini gives
\(\int_Y S\,d\nu=\sum_n\int_Y f_n\,d\nu\).
The numerical series is absolutely convergent because
\[
 \sum_n\left|\int_Yf_n\,d\nu\right|
       \leq\sum_n\int_Y|f_n|\,d\nu<\infty.
\]
Every \(f_n\) is individually integrable, since each term in the
assumed nonnegative finite sum is finite. Thus each signed integral
appearing in this argument was defined before the interchange.
\end{proof}
\end{annexitem}

\begin{annexitem}{annex:p15.layer-cake}{thm:p15-layer-cake}{A measurable subgraph and its section lengths}
\begin{proof}
The second factor is \([0,\infty)\) with the trace Borel
sigma-algebra and restricted Borel Lebesgue measure; it is sigma-finite
by its bounded intervals. The set \(D_g\) has the countable union
description
\[
 D_g=\bigcup_{q\in\mathbb Q,\ q>0}
            \{x:g(x)>q\}\times[0,q).
\]
Indeed, if \(0\leq t<g(x)\), there is a positive rational
\(q\) strictly between \(t\) and \(g(x)\), also when
\(g(x)=\infty\). Conversely, membership in a displayed rectangle
gives \(t<q<g(x)\). Every rectangle is measurable, proving
product measurability. At a fixed \(x\), the section is
\([0,g(x))\) when \(g(x)<\infty\), including the empty set when
\(g(x)=0\), and is \([0,\infty)\) when \(g(x)=\infty\).
Its Lebesgue measure is therefore exactly \(g(x)\) in every case.
At a fixed \(t\geq0\), the other section is \(\{x:g(x)>t\}\).
The section-measure lemma, equivalently Tonelli for \(\mathbf1_{D_g}\),
now gives
\[
 \int_Xg\,d\mu
   =(\mu\times\lambda)(D_g)
   =\int_{[0,\infty)}\mu(\{x:g(x)>t\})\,d\lambda(t).
\]
This also proves measurability of the tail-measure function; no density
or continuity assumption on the distribution of \(g\) was needed.
For a probability space, \(\mu\) is finite and hence sigma-finite.
Putting \(g=Z\) and using the definition of expectation gives the
stated probability tail formula. Including or excluding the endpoint
zero in the outer Lebesgue domain leaves its integral unchanged,
because that singleton is null and the integrand is nonnegative.
\end{proof}
\end{annexitem}
